\section{Exact functors from one Galois category into another}
\label{V.6}

\begin{proposition}
\label{V.6.1}
Let $\cal{C}$, $\cal{C'}$ be two Galois categories,
$H\colon \cal{C}\to \cal{C'}$ a covariant functor, $F'$ a fundamental
functor on~$\cal{C'}$ and $F=F' \circ H$. The following conditions are
equivalent:
\begin{enumerate}
\item[(i)] $H$ is \emph{exact}, \ie left exact and right
exact.
\item[(ii)] $H$ is left exact, transforms finite sums into
finite sums, and epimorphisms into epimorphisms (or again:
objects $\not\approx \emptyset_{\cal{C}}$ into objects $\not\approx
\emptyset_{\cal{C'}}$).
\item[(iii)] $F$ is a fundamental functor on~$\cal{C}$.
\end{enumerate}
\end{proposition}

The implication (i)$\To$(ii) is a general fact about
categories. Moreover the first form given to (ii)
implies the second, as one sees by noting that if $X$ is an object
of~$\cal{C}$, then $X$ is $\not\approx \emptyset_{\cal{C}}$ if and only if the
morphism $X\to e_{\cal{C}}$ is an epimorphism; one will note that
$F$, being assumed left exact, transforms $e_{\cal{C}}$
into~$e_{\cal{C'}}$. The second form given to (ii)
implies~(iii), since $F$, being left exact and hence
pro-representable, falls under~\Ref{V.5.6}
criterion~(iii). Finally (iii) implies~(i), as follows from the
fact that $F$ is exact, and ``conservative'' (\ie satisfies axiom
(G~6) of~\no \Ref{V.4}).

Let then $\Gamma$ be the fundamental groupoid of~$\cal{C}$, and
$\Gamma'$ that of~$\cal{C'}$. Thus if
% original p. 135
$H$ is exact, then $F' \mto F' \circ H$ is a functor from the
groupoid~$\Gamma'$ into the groupoid~$\Gamma$, which we shall
denote by $\tH$:
$$
\tH (F')(X)=F' (H(X))
$$
which can also be written, with the notation $F(X)=E_X (F)$
introduced in~\no \Ref{V.6}:
$$
E_{H(X)} (F')=E_X (\tH(F')).
$$
This last formula shows, taking into account~\Ref{V.5.8}
or~\Ref{V.4.1}, that the exact functor $H$ is determined (up to
unique isomorphism) when one knows the
corresponding functor~$\tH$. Let us fix an~$F'$, and let $F=\tH (F')$;
then $\tH$ defines a homomorphism from~$\pi_{F'} =\Aut(F')$
into $\pi_F =\Aut(F)$:
$$
\tH\colon \pi_{F'} \to \pi_F \quad (F=\tH (F')=F' \circ H).
$$
Moreover the formula above shows (taking into account~\Ref{V.5.8})
that this homomorphism has the following property: for every
finite set $E$ on which $\pi_F$ operates continuously, the group
$\pi_{F'}$ also operates \emph{continuously} by means of
the preceding homomorphism $\pi_{F'} \to
\pi_F$. \hbox{Applying} this to the quotients of~$\pi_F$ by its open
invariant subgroups, one sees that the preceding condition
also means that the homomorphism under consideration is
continuous. Conversely,
let us be given
an object $F$ of~$\Gamma$, an
object $F'$ of~$\Gamma'$ and a continuous homomorphism
$$
u\colon \pi_{F'} \to \pi_F ,
$$
to it there therefore corresponds a functor from~$\cal{C} (\pi)$
into~$\cal{C}(\pi')$, manifestly exact, hence by virtue
of~\Ref{V.4.1} there corresponds to it a functor $H$ from~$\cal{C}$ into
$\cal{C'}$ which is exact, and such that $\tH\colon \pi_{F'} \to \pi_F$
is precisely~$u$. One can also, instead of a
homomorphism of groups, start from a \emph{functor}
$$
U\colon \Gamma' \to \Gamma
$$
which is such that for \emph{every} $F' \in \Gamma'$ (or for \emph{one} $F'
\in \Gamma'$, which amounts to the same) the corresponding homomorphism
$\pi_{F'} \to \pi_F$ is continuous: such a functor is isomorphic to
a functor of the form~$\tH$, where $H\colon \cal{C} \to
\cal{C'}$ is an exact functor determined up to unique
isomorphism. Thus:

\begin{corollary}
\label{V.6.2}
For
% original p. 136
a functor $H\colon \cal{C} \to \cal{C'}$ between Galois categories
to be exact, it is necessary and sufficient that there exist
equivalences $\cal{C}(\pi)\to\cal{C}$ and
$\cal{C'}\to\cal{C}(\pi')$
which transform the functor $H$ into the functor
$\cal{C}(\pi)\to\cal{C}(\pi')$ associated with a homomorphism of
topological groups $\pi'\to\pi$.
\end{corollary}

\begin{corollary}
\label{V.6.3}
Let $\cal{C}$, $\cal{C'}$ be two Galois categories, and
$\Gamma$, $\Gamma'$ their fundamental groupoids. Then the
category of exact functors from~$\cal{C}$ into $\cal{C'}$ is
equivalent to the category of functors $U\colon \Gamma'
\to \Gamma$ having the following property: for every $F'$ in~$\Gamma'$ (or for \emph{one} $F'$ in~$\Gamma'$, which amounts to the
same), setting $F=U(F')$, the homomorphism
$$
\pi_{F'} =\Aut (F')\to \pi_F =\Aut (F)
$$
defined by $U$ is continuous.
\end{corollary}

Consider the fundamental pro-group $\Pi$ of~$\cal{C}$; then an
exact functor $H$ transforms it into a pro-group $H(\Pi)$
of~$\cal{C'}$; we are going to define a homomorphism
$$
\Pi' \to H(\Pi)
$$
(where $\Pi'$ is the fundamental pro-group of~$\cal{C'}$), by the
condition that for every object $F'$ of~$\Gamma'$, the corresponding
homomorphism
$$
F' (\Pi')=\pi_{F'} \to F' (H(\Pi))=\pi_F \qquad
\text{(where $F=F'\circ H=\tH(F')$)}
$$
be the natural homomorphism
$$
\Aut(F')\to \Aut(F' \circ H).
$$
Since the latter is functorial in~$F'$, it indeed defines, by virtue
of~\Ref{V.5.8}, a homomorphism of pro-objects, and in fact of
pro-groups, of~$\cal{C'}$. This homomorphism is said to be
\emph{associated} with the functor~$H$.

Now let $H'$ be a second exact functor, from the
Galois category $\cal{C'}$ into a
Galois category~$\cal{C''}$. It is trivial that one has
$$
{}^t (H' H)=\tH\,\tH'
$$
(N.B. we have here an identity of functors, and not merely a
canonical isomorphism). One has an analogous transitivity
property for the associated homomorphisms of the
fundamental pro-groups.

We
% original p. 137
shall now interpret the properties of the exact
functor $H$ in terms of the corresponding homomorphism
$$
u\colon \pi_{F'} \to \pi_{F} \qquad \text{(where $F' =F'\circ H$)}.
$$
It is convenient to introduce the notion of a \emph{pointed object}
of the Galois category $\cal{C}$ (endowed with its fundamental
functor $F$): it is by definition an object $X$ of~$\cal{C}$
endowed with an element $a$ of~$F(X)$. It is therefore interpreted
as a finite set on which $\pi_F$ operates continuously on the
left, endowed with a point~$a$. Hence the \emph{connected} pointed
objects of~$\cal{C}$ are identified, by virtue of~\Ref{V.5.3}, with the
open subgroups of~$\pi_F$. If $U$ and $V$ are two such
subgroups, corresponding to connected pointed objects
$X$, $Y$ of~$\cal{C}$, then there exists a pointed homomorphism from
$X$ into $Y$ if and only if $U\subset V$, and this homomorphism is
then unique. Of course, the functor $H$ transforms pointed
objects into pointed objects (since $F=F' \circ H$). Let us note
on the other hand that a closed subgroup of a group such as $\pi_F$
is the intersection of the open subgroups that contain it;
consequently,
if
$M$, $N$ are two closed subgroups, then $M\subset N$ if and
only if every open subgroup that contains $N$ also
contains~$M$. Thanks to these remarks, one easily proves
the results that follow:

\begin{proposition}
\label{V.6.4}
Let $X$ be a connected pointed object of~$\cal{C}$, associated with
an open subgroup $U$ of~$\pi_F$. In order that
$U$
contain $u(\pi_{F'})$ (\resp the closed invariant subgroup
generated by~$u(\pi_{F'})$), it is necessary and sufficient that $H(X)$
admit a pointed section (\resp be completely
decomposed).
\end{proposition}

We shall call a \emph{section} --- understood: over the final
object --- of an object $X$ of a Galois category~$\cal{C}$ a
morphism from the final object $e_{\cal{C}}$ into~$X$, which amounts
to giving an element $a$ of~$F(X)$ invariant under
$\pi_F$; if $X$ is pointed, one says that one has a \emph{pointed
section} if it is compatible with the pointed structures
on~$X$ and~$e_{\cal{C}}$, \ie if $a$ is precisely the marked
object of~$F(X)$. Such a section is therefore unique, and exists if
and only if the marked object of~$F(X)$ is invariant
under~$\pi_F$. Finally, an object of a Galois category is said to be
\emph{completely decomposed}
if it is isomorphic to a sum of final objects, \ie if $\pi_F$
operates trivially on $F(X)$ --- a condition evidently
stronger than the existence of a pointed section, when $X$ is
pointed. Proposition~\Ref{V.6.4} follows trivially from the
preceding definitions and remarks.

\begin{corollary}
\label{V.6.5}
For
% original p. 138
$u$ to be trivial, it is necessary and sufficient that for every object
$X$ of~$\cal{C}$, $H(X)$ be completely decomposed.
\end{corollary}

\begin{proposition}
\label{V.6.6}
Let $X'$ be a connected pointed object of~${\cal{C}}'$, associated
with an open subgroup $U'$ of~$\pi_{F'}$. In order that $U'$
contain~$\Ker u$, it is necessary and sufficient that there exist a
connected pointed object $X$ of~$\cal{C}$ and a pointed homomorphism
from the pointed connected component $X'_0$ of~$H(X)$ into~$X'$ (hence
that $X$ be isomorphic as a pointed object to a quotient of the
neutral connected component of the inverse image of a pointed object
of~$\cal{C}$). If $u$ is surjective, the preceding condition
is also equivalent to the following: $X$ is isomorphic to
an~$H(X)$, where $X$ is a pointed object of~$\cal{C}$.
\end{proposition}

(One calls \emph{neutral connected component}
of a pointed object $X$ of a Galois category~$\cal{C}$
the unique connected pointed subobject of~$X$; it corresponds to the
orbit under $\pi_F$ of the marked point of~$F(X)$, by virtue
of~\Ref{V.5.3}). Since the fact that $U'$ contains $\Ker u$ does not
depend on the chosen pointing of~$X'$ (because another
pointing amounts to replacing $U$ by a subgroup
conjugate to~$U$), one sees:

\begin{corollary}
\label{V.6.7}
In order that $U'$ contain~$\Ker u$, it is necessary and sufficient that there exist an
\hbox{object~$X$} of~$\cal{C}$ (which one may assume connected) and a morphism
from a connected component of~$H(X)$ into~$X'$. If $u$ is surjective,
this also means that $X'$ is isomorphic to an object of the
form~$H(X)$.
\end{corollary}

\begin{corollary}
\label{V.6.8}
In order that $u$ be injective, it is necessary and sufficient that for every object
$X'$ of~${\cal{C}}'$ there exist an object $X$ of~$\cal{C}$ and a
homomorphism from a connected component of~$H(X)$ into~$X'$.
\end{corollary}

\begin{proposition}
\label{V.6.9}
The following conditions are equivalent:
\begin{enumerate}
\item[(i)] The homomorphism $u\colon \pi_{F'} \to \pi_F$ is surjective.
\item[(ii)] For every connected object $X$ of~$\cal{C}$, $H(X)$ is
connected.
\item[(iii)] The functor $H$ is fully faithful.
\end{enumerate}
\end{proposition}

This last fact means that for two objects~$X$, $Y$ of~$\cal{C}$,
the natural map
$$
\Hom(X,Y) \to \Hom(H(X),H(Y))
$$
is bijective.

\begin{corollary}
\label{V.6.10}
For
% original p. 139
$u$ to be an isomorphism, it is necessary and sufficient that $H$ be an
equivalence of categories, or again that the following two conditions
be satisfied:
\begin{enumerate}
\item[a)] for every connected object $X$ of~$\cal{C}$, $H(X)$ is connected
\item[b)] every object of~$\cal{C}'$ is isomorphic to an object of the
form~$H(X)$.
\end{enumerate}
\end{corollary}

\begin{proposition}
\label{V.6.11}
Let $H\colon \cal{C}\to \cal{C}'$ and $H'\colon\cal{C}\to \cal{C}''$
be exact functors between Galois categories, and $F''$ a
fundamental functor on~$\cal{C}''$; let us set $F'=F''H'$ and $F=F'H$, and
consider the associated homomorphisms
$$
u'\colon\pi_{F''}\to \pi_{F'}\qquad u\colon\pi_{F'}\to \pi_F.
$$
In order that $\Ker u\subset \Im u'$, \ie in order that $uu'$ be
the trivial homomorphism, it is necessary and sufficient that for every object $X$
of~$\cal{C}$, $H'(H(X))$ be completely decomposed. In order
that $\Ker u \supset \Im u' $, it is necessary and sufficient that for every
connected pointed object $X'$ of~$\cal{C}'$ such that $H'(X')$ admits a
pointed section, there exist an object~$X$ of~$\cal{C}$ and a
homomorphism from a connected component of~$H(X)$ into~$X'$.
\end{proposition}

The first assertion follows from the last statement
of~\Ref{V.6.4}. The second follows from the conjunction
of~\Ref{V.6.4} and~\Ref{V.6.6}.

\begin{remark}
\label{V.6.12}
It is not true in general, under the conditions
of~\Ref{V.6.8}, that $X'$ is isomorphic to an object of the form
$H(X)$. One can show that in order that every connected object (hence every
object) of~$\cal{C}'$ be isomorphic to an object of the form $H(X)$,
it is necessary and sufficient that $u$ be an isomorphism of~$\pi_{F'}$ onto a
subgroup of~$\pi_F$ that is a \emph{direct factor}. In practice, however,
one directly constructs a homomorphism $\pi_F\to\pi_{F'}$ right inverse
to~$u$ by means of a suitable exact functor
from~$\cal{C'}$ into~$\cal{C}$.
\end{remark}

\begin{proposition}
\label{V.6.13}
Let $\cal{C}$ be a Galois category endowed with a fundamental
functor $F$, $S$ a connected object of~$\cal{C}$, and $\cal{C}'$ the
category of objects of~$\cal{C}$ over~$S$. Then
$\cal{C'}$ is a Galois category, and the functor $X\mto
H(X)=X\times S$ from~$\cal{C}$ into~$\cal{C}'$ is exact. Let $a\in
F(S)$, and let $F'$ be the functor from~$\cal{C}'$ into the category
of finite sets defined by
$$
F'(X')=\text{inverse image of~$a$ under }F(X')\to F(S).
$$
Then one has an isomorphism $F\cong F'\circ H$, and the
corresponding homomorphism
$$
u:\pi_{F'}\to \pi_F
$$
% original p. 140
is an isomorphism of~$\pi_{F'}$ onto the open subgroup $U$ of
$\pi_{F}$ that is the stabilizer of the marked element $a$ of~$F(X)$.
\end{proposition}

The proof is left to the reader.

\section{The case of preschemes}
\label{V.7}
Let $S$ be a locally noetherian and
\emph{connected} prescheme, and let
$$
a\colon \Spec(\Omega)\to S
$$
be a geometric point of~$S$, with values in an
algebraically closed field $\Omega$. We shall set
$$
\cal{C}=\text{category of \'etale coverings of~$S$}
$$
and for an object $X$ of~$\cal{C}$, \ie an \'etale covering
$X$ of~$S$, we set
$$
F(X)=\text{set of geometric points of~$X$
over~$a$.}
$$
Thus $F$ becomes a functor on~$\cal{C}$ with values in the
category of finite sets. Properties (G~1) to
(G~6) are satisfied: for (G~1), this is contained in the sorites of
I~\Ref{I.4.6}, (G~2) follows from~\Ref{V.3.4}, (G~3) from~\Ref{V.3.5},
(G~4) is trivial by definition, (G~5) follows from~\Ref{V.3.5} and
from the beginning of \No \Ref{V.2}, finally (G~6) is proved in~\Ref{V.3.7}. One
can therefore apply the results of \No \Ref{V.4}, \Ref{V.5}, \Ref{V.6}. This makes it
possible in particular to define a pro-object $P$ of~$\cal{C}$
representing $F$, called the \emph{universal covering of~$S$
at the point~$a$},
and a topological group $\pi=\Aut(F)=\Aut^0 P$, called the
\emph{fundamental group of~$S$ at~$a$},
and denoted~$\pi_1(S,a)$.
The functor $F$ then defines an equivalence of the
category $\cal{C}$ with the category of finite sets
on which $\pi=\pi_1(S,a)$ operates continuously. This
equivalence therefore makes it possible to interpret the usual
operations of projective limits and finite inductive limits on
coverings (products, fiber products, sums, passage to the
quotient, etc.) in terms of the analogous operations in
$\cal{C}(\pi)$, \ie in terms of the obvious operations on
finite sets on which $\pi$ operates. Let us note moreover, since
the topological connected components of an \'etale covering
are also \'etale coverings, that \emph{an \hbox{object $X$} of~$\cal{C}$ is connected in $\cal{C}$ if and only if it is
topologically connected}; by virtue of~\Ref{V.5.3} this therefore means
that $\pi_1$ operates transitively on~$F(X)$.
% original p. 141
Let us note that in order for an object $X$ of~$\cal{C}$ to be faithfully flat
and quasi-compact over~$S$ (as it is already flat and
quasi-compact over~$S$), it is necessary and sufficient that $X\to S$ be surjective,
\ie be an epimorphism in $\cal{C}$, or again that $X\neq
\emptyset$. One then concludes from criterion \Ref{V.2.6}~(iii) that~$X$
\emph{is a principal covering of
$S$
with group $G$ if and only if it is a principal homogeneous space
under $G$ in the category $\cal{C}$}, (as it was
defined in \No \Ref{V.5}).

If $a'$ is another geometric point of~$S$ (corresponding
to an algebraically closed field $\Omega'$, which may be
different from~$\Omega$ and which may even have a
different characteristic), it defines a fiber functor
$F'=F_{a'}$ from~$\cal{C}$ into the category of finite sets,
which is again exact, hence isomorphic to $F=F_a$. Consequently, the
fundamental groups $\pi_1(S;a)$
for variable $a$ are isomorphic to one another. If one denotes by
$\pi_1(S;a,a')$
the set of isomorphisms (or, what amounts to the same,
the set of homomorphisms) $F_a\to F_{a'}$ of the associated fiber functors,
one thus obtains a \emph{groupoid} whose
set of objects is the set of geometric points
of~$S$, the fundamental groups being the automorphism
groups of the objects of the said groupoid. The set
$\pi_1(S;a',a)$ may be called the \emph{set of classes
of paths from~$a$ to $a'$}. These classes therefore compose in the obvious way. Finally, one can define a pro-group $\Pi_1^S$
of~$\cal{C}$, which one may call the \emph{fundamental pro-group
of}~$S$
or \emph{local system of fundamental groups on}~$S$,
defined up to unique isomorphism by the condition that one
have an isomorphism, functorial in the geometric point $a$ of
$S$:
$$
F_a(\Pi_1^S)=\pi_1(S;a)
$$
(\cf remark~\Ref{V.5.10}). In particular, if $s$ is an ordinary
point of~$S$, the fiber of~$G$ at~$s$ is a pro-group over~$\kres(s)$,
a projective limit of finite \'etale groups over~$\kres(s)$; one could
call this pro-group \emph{the fundamental group of~$S$ at the
ordinary point $s$ of~$S$}, and denote it $\pi_1(S,s)$. By definition, these
points with values in an algebraically closed extension
$\Omega$ of~$\kres(s)$ are the elements of~$\pi_1(S;a)$, where
$a$ is the geometric point of~$S$ defined by the said
extension. In particular, (taking for $S$ the spectrum of a field)
to every field $k$ there is associated canonically and functorially
a pro-group over~$k$, which one could denote $\pi_1(k)$, a
projective limit of finite \'etale groups over~$k$, and whose points
in an algebraically closed extension $\Omega$ of~$k$
are identified
% original p. 142
with the elements of the topological Galois group of~$\bar{k}/k$,
where $\bar{k}$ is the Galois closure of~$k$ in $\Omega$
(\cf \Ref{V.8.1}). This group $\pi_1(k)$ does not yet seem to have
attracted the attention of algebraists.

Now let
$$
f\colon S'\to S
$$
be a morphism from one connected locally noetherian prescheme
into another, let $a'$ be a geometric point of~$S'$ and let
$a=f(a')$ be its direct image in $S$. Then the ``inverse
image'' functor induces a functor from the category $\cal{C}(S)$
of \'etale coverings of~$S$ into the category
$\cal{C}(S')$ of \'etale coverings of~$S'$:
$$
f^\bbullet:\cal{C}(S)\to \cal{C}(S')
$$
One has moreover an isomorphism of functors
$$
F_a\cong F_{a'}\circ f^\bbullet,
$$
so that $f^\bbullet$ is an \emph{exact} functor, to which
the results of \No \Ref{V.6} apply. In particular one has a
canonical homomorphism
$$
u=\pi_1(f;a'):\pi_1(S',a')\to \pi_1(S;a) \qquad (a'=f(a))
$$
which makes it possible to reconstruct the inverse image functor, as an
operation of restriction of the groups of operators. The
properties of the functor $f^\bbullet$ are therefore expressed in a
simple way by the properties of the associated homomorphism of
groups, as was made explicit in
\No \Ref{V.6}. If in particular $S'$ is an \'etale covering of~$S$,
then the homomorphism $u$ is an isomorphism of~$\pi_1(S',a')$ onto the
open subgroup of~$\pi_1(S,a)$ that defines the pointed connected
\'etale covering $S'$ of~$S$ (\ie the stabilizer~$U$
of~$a'\in F_a(S')$ in $\pi_1(S;A)$).

If one wishes to interpret the homomorphisms $\pi_1(f;a')$ for
a variable geometric point $a'$, one must, in accordance
with what was said in \No \Ref{V.6}, consider a
homomorphism
$$
\Pi_1(f):\Pi_1^{S'}\to f^\bbullet(\Pi_1^S)
$$
of pro-groups over~$S$, and take the corresponding homomorphism for
the geometric fibers.

\section{The case of a normal base prescheme}
\label{V.8}
% original p. 143
\begin{proposition}
\label{V.8.1}
Let $S$ be the spectrum of a field $k$, and let $\Omega$ be an
algebraically closed extension of~$k$, defining a geometric
point $a$ of~$S$ with values in $\Omega$. Let
$\bar{k}$ be the separable closure of~$k$ in $\Omega$. Then there
exists a canonical isomorphism of~$\pi_1(S,a)$ onto the topological
Galois group of~$\bar{k}/k$.
\end{proposition}

Let $k'$ be the algebraic closure of~$k$ in $\Omega$; there
therefore corresponds to it a geometric point~$b$ of~$S$, with
values in $k'$. The natural homomorphism of functors $F_b\to F_a$
is evidently an isomorphism, since a $k$-homomorphism of a
finite separable extension of~$k$ into $\Omega$
necessarily takes its values in $\bar{k}$ and a fortiori in
$k'$. On the other hand, the group $\pi'$ of $k$-automorphisms of~$k'/k$
operates in the obvious way on~$F_b$, whence a
homomorphism $\pi'\to \Aut(F_b)\isomto \Aut(F_a)=\pi_1(S;a)$.
On the other hand, it is well known that the natural homomorphism from~$\pi'$
into the group $\pi$ of automorphisms of~$\bar{k}/k$ is an
isomorphism. One thus obtains a canonical homomorphism $\pi\to
\pi_1(S;a)$; it remains to show that it is an isomorphism. Indeed,
this homomorphism is injective, since an element of the kernel is an
automorphism of~$\bar{k}/k$ which induces the identity on every
finite separable subextension, hence is trivial. This
homomorphism is surjective, since if $X$ is a \emph{connected}
\'etale covering of~$S$, hence defined by a finite separable
\emph{extension}~$L/k$, then $\pi$ is transitive
on the set of $k$-homomorphisms of~$L$ into $k'$, as is well
known.

\begin{proposition}
\label{V.8.2}
Let $S$ be a connected, locally noetherian and
normal prescheme, $K=\kres(s)$ its function field $=$ the residue field at
its generic point $s$, and $\Omega$ an
algebraically closed extension of~$K$, defining a geometric
point $a'$ of~$S'=\Spec(K)$ and a geometric
point $a$ of~$S$. Then the homomorphism
$\pi_1(S';a')\to\pi_1(S;a)$ is surjective. When one identifies the
first group with the Galois group of the separable closure
$\bar{K}$ of~$K$ in~$\Omega$ (\cf \Ref{V.8.1}), then the kernel of
the preceding homomorphism corresponds by Galois theory
to the subextension of~$\bar{K}/K$ composed of the
finite extensions of~$K$ in $\Omega$ which are unramified over~$S$.
\end{proposition}

The first assertion means that the inverse image on~$S'$ of a
connected \'etale covering $X$ of~$S$ is connected, \ie that $X$
is integral; this is none other than (I~\Ref{I.10.1}). The kernel of
the preceding homomorphism is then interpreted as
consisting of the automorphisms of~$\bar{K}/K$ which induce
the identity on the sets~$F_a(X)$,
% original p. 144
where one may assume the \'etale covering $X$ of~$S$
connected. But this means that this automorphism induces
the identity on the finite subextensions of~$\bar{K}/K$ which are
unramified over~$S$, which proves the last assertion.

\begin{remarkstar}
Thanks to this interpretation of the fundamental group of the
normal prescheme $S$ in terms of the usual Galois theory,
the definition had been known in this case for a
long time.
\end{remarkstar}

\section{The case of non-connected preschemes: multi-Galois
categories}
\label{V.9}

Let $S$ be a locally noetherian prescheme, and let
$(S_i)_{i\in I}$ be its connected components. Then the category
$\cal{C}(S)$ of \'etale coverings of~$S$ is equivalent
to the product category of the $\cal{C}(S_i)$, which
are interpreted in terms of the fundamental groups of the $S_i$, once
a geometric point has been chosen in each $S_i$. In the
application of descent theory for \'etale morphisms, it is
sometimes inconvenient to choose, for every $S_i$,
a geometric point of~$S_i$. It is then more convenient
to resort to the natural generalization of~\Ref{V.5.8}
in order to interpret $\cal{C}(S)$ as a category of functors
on the groupoid of geometric points of~$S$,
considered as the sum of the groupoids corresponding to the
connected components of~$S$; the functors in question are the
functors with values in the category of finite sets
satisfying the continuity property analogous to
the one invoked in~\Ref{V.5.8}. In practice, one will have a family
$(a_t)_{t\in E}$ of geometric points of~$S$ such that
every connected component $S_i$ of~$S$ contains at least one of them, and one
can then, as in~\Ref{V.5.10}, replace the groupoid of
all geometric points of~$S$ by the
analogous groupoid whose underlying set is~$E$. Of course, these
considerations ought to fit into
general definitions concerning the categories which are
equivalent to product categories of categories
of the form $\cal{C}(\pi)$, and which one may call
\emph{multi-Galois categories}.
We shall leave the details to the reader.
